\documentclass{report}
\usepackage{amsmath,amssymb}
\setlength{\parindent}{0mm}
\setlength{\parskip}{1em}
\begin{document}
\begin{center}
\rule{6in}{1pt} \
{\large Piyush K Sao \\
{\bf Scalable Sparse Direct Solver for Hybrid Architecture}}

RM\#1343 \\ 266 ferst Dr NW \\ Atlanta GA 30318
\\
{\tt piyush3@gatech.edu}\\
Xiaoye S.  Li \\
Richard W. Vuduc\end{center}

We present new algorithmic techniques for improving the scalability
and efficiency of sparse direct solvers, specifically for distributed
memory systems comprised of hybrid multicore CPU systems and hardware
co-processors. This work extends the algorithm in SUPERLU\_DIST, which is
both right-looking and statically pivoted.

First, we present a novel algorithm for exploiting intra-node
co-processors, which we call the HALO (Highly Asynchronous Lazy Offload)
algorithm. HALO combines highly aggressive use of asynchrony with
accelerated offload, lazy updates, and data shadowing (a la halo or ghost
zones). These techniques hide and reduce communication, whether to local
memory, across the network, or over PCIe. We evaluated our implementation
on both NVIDIA GPUs and Intel Xeon-Phi accelerated systems. We use
various realistic test problems for both single-node and multi-node
configurations, finding that our implementation achieves
speedups of up to 2.5x over an already efficient multicore CPU implementation.

Secondly, we present new algorithmic techniques to improve the
scalability of sparse direct solvers at high core counts. Our techniques
aggressively exploit elimination tree parallelism and replicate data to
reduce per process communication volume. The net effect is to overlap
communication with computation. We give preliminary experimental results
that show the efficacy of these techniques.

Beyond sparse direct solvers, we believe our proposed techniques can
apply to wide range of other sparse linear algebra computations.


\end{document}
